\clearpage
\VerdierPhysicalPageAnchor{0228}%
\begingroup\small\setlength{\parskip}{0pt}\setlength{\emergencystretch}{1.5em}
\setlength{\abovedisplayskip}{4pt plus 2pt minus 2pt}%
\setlength{\belowdisplayskip}{4pt plus 2pt minus 2pt}%

\VerdierSemanticAnchor{ch4d-4-1-4}
\medskip\noindent\textbf{4.1.4.}\quad
由 (4.1.3)，对 $\operatorname{comp}(\mathcal A)$ 的每个对象 $Y$，

\VerdierSemanticAnchor{ch4d-4-1-4-spec-k}
\[
\operatorname{Spec}_1(Y)=\bigl(Y(p,q),\delta(p,q,r)\bigr),
\]
由 (4.1.2.3)、(4.1.2.4) 与 (4.1.2.6) 定义，是取值于
$K(\mathcal A)$ 的谱对象。谱对象 $\operatorname{Spec}_1(Y)$ 函子性地依赖于
$\operatorname{comp}(\mathcal A)$ 的对象 $Y$。因此得到一个函子

\VerdierSemanticAnchor{ch4d-4-1-4-functor}
\[
Y\longmapsto\operatorname{Spec}_1(Y)
\]
从 $\mathcal A$ 的复形范畴到取值于 $K(\mathcal A)$ 的谱对象范畴。这个函子
称为\emph{取值于 $K(\mathcal A)$ 的第一典范谱对象函子}。对
$\operatorname{comp}(\mathcal A)$ 的每个对象 $Y$，谱对象
$\operatorname{Spec}_1(Y)$ 称为\emph{与复形 $Y$ 相联系的、取值于
$K(\mathcal A)$ 的第一谱对象}。应注意，当 $Y$ 是
$\operatorname{comp}^{*}(\mathcal A)$ 的对象，且 $*$ 取 $+$、$-$、$b$ 或
“空”时，与复形 $Y$ 相联系的、取值于 $K(\mathcal A)$ 的第一谱对象实际上
取值于 $K^{*}(\mathcal A)$。

\VerdierSemanticAnchor{ch4d-4-1-5}
\medskip\noindent\textbf{4.1.5.}\quad
设 $\mathcal A$ 是阿贝尔范畴，$Y$ 是 $\mathcal A$ 的一个复形。与复形
$Y$ 相联系的第一谱对象在典范函子
$K(\mathcal A)\longrightarrow D(\mathcal A)$ 下的像，是一个取值于
$D(\mathcal A)$ 的谱对象，称为\emph{与复形 $Y$ 相联系的、取值于
$D(\mathcal A)$ 的第一谱对象}。通常，我们简称为\emph{与复形 $Y$ 相联系的
第一谱对象}，并（滥用记号）写作

\VerdierSemanticAnchor{ch4d-4-1-5-spec-d}
\[
\operatorname{Spec}_1(Y)=\bigl(Y(p,q),\delta(p,q,r)\bigr)
\]
来表示与复形 $Y$ 相联系的、取值于 $D(\mathcal A)$ 的第一谱对象。第一典范
谱对象函子与典范函子 $K(\mathcal A)\longrightarrow D(\mathcal A)$（延拓到
谱对象）的复合称为\emph{取值于 $D(\mathcal A)$ 的第一典范谱对象函子}。
需要强调，这是从 $\mathcal A$ 的复形范畴到 $D(\mathcal A)$ 的谱对象范畴的
函子，并且\emph{一般而言，这个函子不通过复形同伦范畴分解}。

\VerdierSemanticAnchor{ch4d-4-1-6}
\medskip\noindent\textbf{4.1.6.}\quad
$\mathcal A$ 的对象所成的 $n$-重复形可以构造带有滤过的单复形，从而构造
取值于 $K(\mathcal A)$ 的谱对象。设 $X$ 是一个 $n$-重复形（第一章，
2.1.2），并设 $m$ 是整数，$1\leq m\leq n$。对每个
$\sigma\in\mathbb Z[n]$（记号见（第一章，2.1.1）），以 $\sigma_m$ 表示其
第 $m$ 个坐标。

\VerdierSemanticAnchor{ch4d-4-1-6-family-setup}
首先定义一族复形以及复形态射。对
$\mathcal{F}\ell_1(\widetilde{\mathbb Z})$ 的每一对 $(p,q)$，令
${}^{m}X(p,q)$ 为
\endgroup
